% Generated draft: requires mathematical and source review.
\documentclass[11pt]{article}
\usepackage[margin=25mm]{geometry}
\usepackage{fontspec}
\setmainfont{Latin Modern Roman}
\usepackage{amsmath,amssymb,mathtools,unicode-math}
\setmathfont{Latin Modern Math}
\usepackage{enumitem,needspace,xurl,hyperref}
\setlength{\parindent}{0pt}
\setlength{\parskip}{5pt}
\setlength{\emergencystretch}{3em}
\newcommand{\N}{\mathbb{N}}
\newcommand{\Z}{\mathbb{Z}}
\newcommand{\Q}{\mathbb{Q}}
\newcommand{\R}{\mathbb{R}}
\newcommand{\C}{\mathbb{C}}
\newcommand{\F}{\mathbb{F}}
\newcommand{\A}{\mathbb{A}}
\newcommand{\PP}{\mathbb{P}}
\newcommand{\E}{\mathbb{E}}
\newcommand{\eps}{\varepsilon}
\newcommand{\dd}{\,\mathrm{d}}
\begin{document}
\section*{UnsolvedMath ID: 701385}
\textbf{Source problem code:} catalogue entry
\subsection{Definitions and mathematical statement}
For a group \(X\), define its lower central series by
\[
\gamma_1(X)=X,\qquad \gamma_{n+1}(X)=[\gamma_n(X),X]\quad(n\geq 1),
\]
where \([U,V]\) denotes the subgroup generated by the commutators \(u^{-1}v^{-1}uv\) with \(u\in U\) and \(v\in V\). The group \(X\) is \emph{residually nilpotent} if
\[
\bigcap_{n\geq 1}\gamma_n(X)=\{1\}.
\]
A \emph{central extension} of a group \(Y\) is a surjective group homomorphism \(p:X\twoheadrightarrow Y\) such that \(\ker p\subseteq Z(X)\), where \(Z(X)\) is the center of \(X\).

Does there exist a group \(G\) satisfying both of the following conditions?
\begin{enumerate}
\item For every group \(E\) and every surjective homomorphism \(p:E\twoheadrightarrow G\),
\[
\ker p\subseteq Z(E)\quad\Longrightarrow\quad
\bigcap_{n\geq 1}\gamma_n(E)=\{1\}.
\]
\item There exist groups \(H,K\) and surjective homomorphisms
\[
K\xrightarrow{\ q\ }H\xrightarrow{\ r\ }G
\]
such that
\[
\ker q\subseteq Z(K),\qquad \ker r\subseteq Z(H),\qquad
\bigcap_{n\geq 1}\gamma_n(K)\neq\{1\}.
\]
\end{enumerate}
All groups and central kernels are unrestricted; the supplied statement imposes no finiteness, finite generation, or splitting hypotheses. Centrality is required for each of the two successive extensions separately.
\subsection{Short English statement}
Does there exist a group whose every central extension is residually nilpotent, but which admits two successive central extensions whose final group is not residually nilpotent? The supplied registry labels this problem open; that status has not been independently verified.
\subsection{Sources}
\begin{itemize}
\item Local UnsolvedMath registry, permanent ID \(701385\), exact source problem code \texttt{catalogue entry}, title \emph{Kourovka Notebook Problem 17.83}. The supplied record attributes the problem to R. Mikhailov, proposed in 2010, and labels it open.
\item Supplied source attribution: \emph{The Kourovka Notebook (Unsolved Problems in Group Theory)}, 17th issue (2010), Problem 17.83; also the 21st issue edition (2026), arXiv:1401.0300v48, \url{https://arxiv.org/abs/1401.0300v48}. The supplied collection description names E. I. Khukhro and V. D. Mazurov as editors. These bibliographic details are supplied attribution and have not been independently verified.
\end{itemize}
\end{document}
